
\subsubsection{2.1 Weight-Space Property Prediction}

Unterthiner et al. [2020] demonstrated that test accuracy can be predicted from weight tensors using flat MLP encoders that concatenate all weight tensors into a single vector. They report Spearman ρ ≈ 0.5–0.7 on a proprietary model zoo and establish Spearman rank correlation as the standard evaluation metric for this task.

Eilertsen et al. [2020] extend weight-space property prediction to training objective classification on CNN weights. Both works use flat concatenation without exploiting permutation symmetry.

\subsubsection{2.2 Symmetry-Aware Encoders}

Navon et al. [2023] introduce Neural Functional Networks (NFNs): equivariant weight-space architectures that respect the neuron-permutation symmetry of feedforward networks. Zhou et al. [2023] independently develop Neural Functional Transformers (NFTs). Kofinas et al. [2024] propose graph neural networks with equivariant message passing for weight spaces. All three works report improved performance over flat MLP baselines but do not match encoder capacities or report bootstrap CIs on Δρ.

Zaheer et al. [2017] prove that any permutation-invariant function on sets decomposes as ρ(Σᵢφ(xᵢ)) (Deep Sets architecture). The Deep Sets architecture has not been benchmarked as an intermediate baseline on model zoo accuracy prediction tasks.

\subsubsection{2.3 Model Zoo Benchmarks}

Schurholt et al. [2022] release the MNIST-CNN and CIFAR-10 model zoos for weight-space learning research. Schurholt et al. [2023] compare multiple encoder architectures on these zoos without capacity matching, bootstrap CIs, or a Deep Sets baseline.

\subsubsection{2.4 Permutation Symmetry}

The permutation symmetry structure of feedforward networks is characterized by Navon et al. [2023] and Zhou et al. [2023]. Related work on linear mode connectivity \citep{Ainsworthetal2023} and loss landscape symmetry \citep{Entezarietal2022} addresses permutation symmetry in the context of model merging. This work focuses on the encoder design implication: structural equivariance avoids the need to learn invariance from data.

